\documentclass[12pt,a4paper]{article}

% 中文支持
\usepackage[UTF8]{ctex}

% 数学公式
\usepackage{amsmath}

% 图片
\usepackage{graphicx}

% 超链接
\usepackage{hyperref}

% 页面设置
\usepackage{geometry}
\geometry{a4paper,scale=0.8}

\title{
Tree-Ring Watermarks: Invisible Fingerprints for Diffusion Images\\
基于扩散模型初始噪声空间的生成图像水印技术研究
}

\author{张志杰}

\date{\today}


\begin{document}

\maketitle


\begin{abstract}

随着扩散模型的发展，人工智能生成图像的质量不断提升，
如何实现生成内容的来源追踪和版权保护成为重要研究问题。
本文总结论文
Tree-Ring Watermarks: Invisible Fingerprints for Diffusion Images
提出的一种基于扩散模型初始噪声空间的不可见水印方法。
该方法通过在Fourier空间中修改初始噪声，
将隐藏指纹嵌入扩散模型生成过程，并利用DDIM inversion实现水印检测。
本文进一步分析相关生成模型水印技术的发展现状，
并讨论未来研究方向。

\end{abstract}


\section{研究背景与意义}

随着人工智能生成内容（Artificial Intelligence Generated Content, AIGC）
技术的快速发展，基于深度学习的生成模型已经成为人工智能领域的重要研究方向。
近年来，以扩散模型（Diffusion Model）为代表的生成模型在文本到图像生成任务中取得了显著进展。
例如，Stable Diffusion、DALL-E 等模型能够根据自然语言描述生成具有较高真实性和艺术表现力的图像，
推动了人工智能生成内容在数字艺术、辅助设计以及内容创作等领域的发展。


然而，生成模型能力的快速提升也带来了新的安全问题。
由于当前生成图像的真实性不断提高，普通用户已经难以仅凭视觉判断图像是否由人工智能生成。
这可能导致虚假信息传播、身份伪造、版权归属困难等问题。
因此，如何实现人工智能生成内容的来源追踪和可信认证成为当前的重要研究方向。


数字水印技术为解决上述问题提供了一种有效方法。
传统数字水印通常通过修改图像像素或者频域信息，将隐藏信息嵌入图像中。
然而，对于扩散模型生成图像而言，直接修改最终输出图像容易受到裁剪、压缩以及图像编辑等攻击。
同时，传统方法没有充分利用扩散模型自身的生成机制。


因此，如何在保证生成图像质量的同时，
将水印信息嵌入扩散模型生成过程，
实现不可见、鲁棒且可检测的生成内容指纹，
成为当前人工智能安全领域的重要研究问题。

\section{国内外研究现状}

随着人工智能生成模型的发展，针对生成内容的水印与溯源技术逐渐成为研究热点。
相关研究主要经历了从传统数字水印、生成模型水印到扩散模型水印的发展过程。


\subsection{传统数字水印技术}

传统数字水印技术主要通过修改数字媒体中的像素信息或者频域信息，
将特定信息隐藏到图像中，以实现版权保护和内容认证。
常见方法包括基于空间域的像素修改方法以及基于频域变换的DCT、DWT等方法。

这类方法经过长期发展已经较为成熟，
但其主要针对普通数字图像设计。
当应用于人工智能生成图像时，
由于生成内容可能经历重新生成、编辑以及复杂图像处理，
传统水印方法容易受到攻击，
同时无法充分利用生成模型自身的结构特点。


\subsection{生成模型水印技术}

随着GAN等深度生成模型的发展，
研究者开始关注如何证明生成内容的来源。
相关方法主要包括模型参数水印和生成结果水印两类。

模型参数水印方法通过修改生成模型内部参数，
使模型携带特定身份信息。
然而，该方法主要用于验证模型本身，
难以直接证明某一张具体图像的生成来源。

另一类方法则直接在生成结果中加入隐藏信息。
虽然能够实现图像级别的检测，
但由于水印位于最终图像空间，
容易受到裁剪、压缩以及图像编辑攻击。


\subsection{扩散模型水印技术}

近年来，扩散模型逐渐成为文本到图像生成领域的重要技术。
与传统生成模型不同，
扩散模型通过从随机噪声逐步去噪生成图像，
因此其生成过程提供了新的水印嵌入空间。

已有研究开始探索在latent空间、采样过程以及初始噪声中嵌入水印。
其中，Tree-Ring Watermarks进一步利用扩散模型初始噪声的特点，
将水印从传统图像空间转移到生成过程空间，
实现了更加隐蔽和鲁棒的生成图像指纹。

\section{Tree-Ring Watermarks核心技术}

\subsection{扩散模型生成过程}

扩散模型（Diffusion Model）是一类基于概率生成过程的深度生成模型。
与传统生成模型不同，扩散模型通过模拟逐步加噪和去噪过程实现数据生成。
在图像生成任务中，模型首先从随机噪声开始，
然后通过多步反向扩散过程逐渐恢复出清晰图像。


在扩散模型生成过程中，初始随机噪声是影响最终生成结果的重要因素。
通常情况下，扩散模型从满足高斯分布的噪声 $z_T$ 开始，
经过反向采样过程得到最终图像 $x_0$：

\begin{equation}
z_T \rightarrow z_{T-1} \rightarrow ... \rightarrow z_0
\end{equation}


由于初始噪声决定了整个采样轨迹，
因此对噪声空间进行适当修改可以影响最终生成结果。
Tree-Ring Watermarks正是利用这一特点，
不直接修改生成后的图像，
而是在生成过程的初始噪声中嵌入隐藏水印信息。\subsection{Fourier空间水印嵌入方法}

为了避免直接修改初始噪声导致生成质量下降，
Tree-Ring方法选择在Fourier频域中嵌入水印信息。

首先，对初始噪声 $z_T$ 进行Fourier变换：

\begin{equation}
Z_T=\mathcal{F}(z_T)
\end{equation}

其中，$\mathcal{F}$表示Fourier变换。
通过频域表示，可以更加灵活地控制水印信息的嵌入位置。


作者设计了一种环形（Ring）结构作为水印区域。
通过定义二值掩码矩阵 $M$，
选择Fourier空间中的特定区域用于存放水印信息。
其中，$M=1$表示该区域使用水印替换，
$M=0$表示保持原始噪声信息。

Tree-Ring通过以下方式生成带水印噪声：

\begin{equation}
z_T^{*}
=
\mathcal{F}^{-1}
(
M\odot K+
(1-M)\odot\mathcal{F}(z_T)
)
\end{equation}

其中：

\begin{itemize}
    \item $K$表示预定义的水印密钥（watermark key）；
    \item $M$表示水印区域掩码；
    \item $\odot$表示逐元素乘法；
    \item $\mathcal{F}^{-1}$表示逆Fourier变换。
\end{itemize}

该公式表示，在Fourier空间中，
Tree-Ring仅替换由掩码指定的环形区域，
其余频率信息仍保持原始随机噪声状态。
因此，该方法能够在保持扩散模型正常生成能力的同时，
将隐藏指纹嵌入生成过程。

\subsection{基于DDIM inversion的水印检测}

由于Tree-Ring水印并未直接存在于最终图像像素中，
因此无法通过简单图像分析进行检测。
为了恢复生成过程中的初始噪声，
论文采用DDIM inversion方法进行反向推断。


给定生成图像$x_0$，
DDIM inversion可以近似恢复其对应的初始噪声：

\begin{equation}
x_0 \rightarrow \hat{z}_T
\end{equation}


随后，方法将恢复得到的噪声转换到Fourier空间，
并检测其中是否存在预定义的Ring结构。
通过计算恢复噪声与水印密钥之间的相似程度，
可以判断该图像是否由带水印的扩散模型生成。


Tree-Ring的检测过程可以概括为：

\begin{equation}
Image
\rightarrow
DDIM\ inversion
\rightarrow
\hat{z}_T
\rightarrow
Fourier\ analysis
\rightarrow
Watermark\ detection
\end{equation}



\section{相关并行技术}

随着生成模型的发展，人工智能生成内容的可信认证问题逐渐受到关注。
目前相关技术主要包括传统数字水印、生成模型水印以及基于生成过程的水印方法。
这些方法从不同角度解决生成内容来源追踪问题。


\subsection{传统数字图像水印技术}

传统数字水印技术主要针对数字媒体版权保护问题，
通过将隐藏信息嵌入图像像素空间或者频率空间，
实现内容认证和版权追踪。

常见方法包括基于空间域的水印方法以及基于DCT、DWT等频域变换的方法。
频域水印通过修改图像中的特定频率分量，
能够在保证视觉质量的同时提高一定的鲁棒性。

然而，传统水印方法通常假设水印是在已有图像基础上的附加信息。
对于人工智能生成图像而言，
生成过程本身包含大量随机性，
因此仅修改最终图像可能无法充分利用生成模型结构特点。


\subsection{生成模型水印技术}

随着深度生成模型的发展，
研究者开始探索针对GAN等生成模型的水印方法。

一种思路是模型参数水印（Model Watermarking），
即通过修改模型参数，使生成模型自身携带特定身份信息。
该方法可以用于模型版权保护，
但通常只能证明某个模型的归属，
难以判断具体生成图像的来源。


另一类方法是在生成结果中嵌入水印。
例如，通过训练编码器和解码器，
学习将隐藏信息嵌入生成图像。
这种方法能够实现图像级别检测，
但是水印仍然存在于输出图像空间，
容易受到图像编辑和后处理攻击。


\subsection{扩散模型水印技术}

近年来，扩散模型成为图像生成领域的重要技术路线。
与传统生成模型不同，
扩散模型通过从随机噪声逐步去噪生成图像，
因此提供了新的水印嵌入位置。

部分研究尝试在latent空间或者生成模型参数中加入水印，
例如通过微调模型使其产生带水印的数据。

与这些方法相比，
Tree-Ring Watermarks不需要重新训练模型，
而是直接修改扩散过程中的初始噪声，
在Fourier空间中嵌入结构化水印。
这种方法能够利用扩散模型生成过程的特点，
实现不可见且具有较强鲁棒性的生成内容指纹。 


\section{未来研究方向}

虽然Tree-Ring Watermarks在扩散模型生成图像水印领域取得了较好的效果，
但随着生成模型技术的发展，
仍然存在一些值得进一步研究的问题。


\subsection{提升复杂攻击环境下的鲁棒性}

当前生成图像可能会经历更加复杂的处理过程，
例如大规模图像编辑、重新生成、风格迁移以及多模型转换等。
未来研究需要设计更加稳定的水印表示方式，
使水印能够在复杂攻击环境下仍然保持较高检测准确率。


\subsection{提高跨模型泛化能力}

目前许多生成模型具有不同的结构和采样方式，
例如不同类型的扩散模型以及未来出现的新型生成模型。
如何设计一种能够适用于不同生成架构的通用水印方法，
是人工智能内容认证领域的重要方向。


\subsection{增加水印信息容量}

Tree-Ring方法主要用于判断图像是否由特定扩散模型生成，
即实现生成内容来源认证。
未来可以进一步研究高容量水印技术，
使水印不仅能够表示模型身份，
还可以包含生成时间、用户信息以及版权信息等更加丰富的数据。


\subsection{结合人工智能内容安全体系}

随着AIGC应用规模扩大，
生成内容认证可能不仅依赖单一水印技术，
而需要与内容检测、数字身份认证以及平台监管机制结合。
未来的生成模型安全体系可能需要形成从生成、传播到验证的完整闭环。


\section{总结}

本文对Tree-Ring Watermarks:
Invisible Fingerprints for Diffusion Images进行了总结。

针对人工智能生成图像难以追踪来源的问题，
该方法提出了一种基于扩散模型初始噪声空间的不可见水印方案。
与传统在生成结果中添加水印的方法不同，
Tree-Ring将水印信息嵌入扩散过程的初始噪声，
并通过Fourier空间中的结构化模式实现隐藏指纹。

在检测阶段，
该方法利用DDIM inversion恢复生成过程中的噪声信息，
再通过分析其中的水印结构完成图像来源判断。

总体而言，
Tree-Ring Watermarks展示了一种利用生成模型内部机制进行内容认证的新思路，
为未来人工智能生成内容的版权保护和安全管理提供了重要参考。


\bibliographystyle{plain}

\begin{thebibliography}{99}

\bibitem{treering}
Yuxin Wen et al.
Tree-Ring Watermarks: Invisible Fingerprints for Diffusion Images.
NeurIPS, 2023.

\end{thebibliography}


\end{document}